\documentclass[11pt]{article}
\usepackage[margin=1in]{geometry}
\usepackage{amsmath,amssymb,amsthm,hyperref}
\newtheorem{theorem}{Theorem}
\newcommand{\C}{\mathbb C}
\title{A periodic counterexample to the generated Taylor-polynomial\ zero-set conjecture}
\author{Principia Math In Progress}
\date{11 September 2026}
\begin{document}
\maketitle

\begin{abstract}
Conjecture 47 of Chung asserts a circular zero-set limit for a rationally
generated polynomial sequence when its fixed quadratic factor has nonpositive
discriminant. The admissible choice $a=b=c=1$ refutes the conjecture:
one infinite subsequence has all zeros on the unit circle, whereas a second
has an additional zero at the origin. Consequently the upper and lower
zero-set limits differ, and the limit required by the source does not exist.
\end{abstract}

\section*{Task 1: Resolve MathDB problem 358323}
The problem, Conjecture 47 on printed page 119 of \cite{chung}, concerns
the polynomials defined by
\begin{equation}\label{eq:gf}
 \sum_{m\geq0}P_m(z)t^m=
 \frac{1}{(at^2+bt+c)(1-tz)},\qquad
 a,b,c\in\mathbb R\setminus\{0\},\quad b^2-4ac\leq0.
\end{equation}
If $t_1$ is a zero of $at^2+bt+c$ of smallest modulus, the proposed conclusion is
\[
 \lim_{m\to\infty}Z(P_m)=\{z\in\C:|z|=1/|t_1|\}.
\]
Here $Z(P)$ is the set of complex zeros of $P$. Definitions 38--41 on
printed pages 92--93 of the source specify the meaning of the limit.
Writing $B(w,r)=\{z:|z-w|<r\}$, they are the usual set limits
\begin{align*}
 w\in\limsup_m Z(P_m)
 &\iff \text{every }B(w,r)\text{ meets }Z(P_m)
       \text{ for arbitrarily large }m,\\
 w\in\liminf_m Z(P_m)
 &\iff \text{every }B(w,r)\text{ meets }Z(P_m)
       \text{ for all sufficiently large }m.
\end{align*}
The source defines $\lim_m Z(P_m)$ when these two sets are equal.

\begin{theorem}\label{thm:counterexample}
For $a=b=c=1$ in \eqref{eq:gf}, and every integer $r\geq0$,
\begin{align}
 P_{3r+1}(z)&=(z-1)\sum_{j=0}^{r}z^{3j},\label{eq:first}\\
 P_{3r+2}(z)&=z(z-1)\sum_{j=0}^{r}z^{3j}.
 \label{eq:second}
\end{align}
In particular,
\[
 0\in\limsup_m Z(P_m)\setminus\liminf_m Z(P_m).
\]
Thus Conjecture 47 is false even for positive $a,b,c$ and strictly negative
discriminant.
\end{theorem}

\begin{proof}
The parameters are nonzero and $b^2-4ac=-3<0$. The two quadratic zeros
are $e^{2\pi i/3}$ and $e^{-2\pi i/3}$; both have modulus one. Thus the
conjectured circle is $|z|=1$.

Define coefficients $c_k$ by
\[
 \frac1{1+t+t^2}=\frac{1-t}{1-t^3}
 =\sum_{j\geq0}(t^{3j}-t^{3j+1}).
\]
These identities hold as formal power series, so
\[
 c_{3j}=1,\qquad c_{3j+1}=-1,\qquad c_{3j+2}=0.
\]
Multiplication by $(1-tz)^{-1}$ gives
\begin{equation}\label{eq:coeff}
 P_m(z)=\sum_{k=0}^{m}c_k z^{m-k}.
\end{equation}
For $m=3r+1$, grouping the nonzero coefficients in pairs yields
\[
 P_{3r+1}(z)
 =\sum_{j=0}^{r}\bigl(z^{3r+1-3j}-z^{3r-3j}\bigr)
 =(z-1)\sum_{j=0}^{r}z^{3j}.
\]
For $m=3r+2$ the extra coefficient $c_{3r+2}$ is zero,
so \eqref{eq:coeff} gives $P_{3r+2}(z)=zP_{3r+1}(z)$.
This proves both factorization identities without division by a polynomial.

Every zero of $S_r(z)=\sum_{j=0}^{r}z^{3j}$ lies on the unit circle:
if $S_r(z)=0$, the polynomial identity
\[
 (z^3-1)S_r(z)=z^{3r+3}-1
\]
implies $z^{3r+3}=1$ and hence $|z|=1$.
(For $r=0$, $S_r=1$ has no zeros.) Formula \eqref{eq:first}
therefore places every zero of $P_{3r+1}$ on $|z|=1$.
Formula \eqref{eq:second} places a zero of $P_{3r+2}$ exactly at $0$.

Every neighborhood of $0$ consequently meets $Z(P_{3r+2})$ for every $r$.
Thus $0$ belongs to the upper limit. But the fixed ball $B(0,1/2)$
is disjoint from $Z(P_{3r+1})$ for every $r$, so $0$ does not belong
to the lower limit. The required limit does not exist.
Moreover, even if the conjectured limit were interpreted as the upper
limit alone, its claimed circle would still omit the point $0$.
\end{proof}

\paragraph{Scope.}
The theorem resolves the universal conjecture by an exact infinite-family
counterexample. It makes no assertion about weak convergence of normalized
zero-counting measures, which is a different question: a single exceptional
zero has mass tending to zero in such a normalization.

\paragraph{Verification and status.}
The accompanying integer-arithmetic verifier constructs $P_m$ independently
from the cubic denominator in \eqref{eq:gf} and checks the coefficient and
factorization identities. These finite checks supplement the proof for all
$r$. A same-agent dependency audit and dated literature/status audit are
saved alongside this file. MathDB was checked to be open with no posted
claims on 11 September 2026; the searched primary sources disclosed no
existing resolution. No external peer review is claimed.

\section*{Task 2: Classify the root-phase obstruction}
For $b^2-4ac<0$, express the quadratic zeros as $\rho e^{\pm i\theta}$.
Determine which rational values of $\theta/\pi$ force the two zero-set
limits to differ, and whether irrational values can have the same effect.

\section*{Task 3: Determine both limiting zero sets}
Compute the full lower and upper zero-set limits for $a=b=c=1$, including
the subsequence $m\equiv0\pmod3$; then extend the description to arbitrary
nonzero real $a,b,c$ with nonpositive discriminant.

\section*{Task 4: Establish a measure version of the limit}
For the polynomials in \eqref{eq:gf}, determine whether the normalized
zero-counting measures converge weakly to uniform arclength on
$|z|=1/|t_1|$ throughout the nonpositive-discriminant range, with zeros
counted with multiplicity.

\section*{Task 5: Quantify exceptional zeros}
For each fixed admissible quadratic and $\varepsilon>0$, bound the number
of zeros lying outside the annulus
$\bigl||z|-1/|t_1|\bigr|<\varepsilon$ and determine the possible limits
of these exceptional zeros along subsequences.

\begin{thebibliography}{9}
\bibitem{chung} J. Chung, \emph{Zero Distribution of Generated Taylor
Polynomials}, MSc project, California State University, Fresno, May 2022;
\href{https://arxiv.org/abs/2305.05130}{arXiv:2305.05130v1}, 9 May 2023.
Definitions 38--41, pp.~92--93; Conjecture 47, p.~119.
\end{thebibliography}
\end{document}
